\chapter{Le langage SQL (SGBD Oracle)}

Le langage \textbf{SQL} (Structured Query Language) est le langage standard utilisé pour interagir avec une base de données relationnelle. Il se compose de trois parties :

\begin{itemize}
    \item \textbf{LDD} (Langage de Définition des Données) : permet de créer, modifier ou supprimer la structure des tables (schémas).\\
    Exemples : \texttt{CREATE TABLE}, \texttt{DROP TABLE}, \texttt{ALTER TABLE}.
    \item \textbf{LMD} (Langage de Manipulation des Données) : permet d’ajouter, supprimer, modifier et interroger les données.\\
    Exemples : \texttt{INSERT}, \texttt{DELETE}, \texttt{UPDATE}, \texttt{SELECT}, \texttt{COMMIT}/\texttt{ROLLBACK}.
    \item \textbf{LCD} (Langage de Contrôle des Données) : utilisé pour gérer les droits d’accès.\\
    Exemples : \texttt{GRANT}, \texttt{REVOKE}.
\end{itemize}

\section{Introduction au SGBD Oracle}

Oracle est un \textbf{système de gestion de bases de données relationnelles} (SGBDR), développé par Oracle Corporation depuis 1981.

La version de référence du cours est \textbf{Oracle 12c} (sortie en 2013), suivie par Oracle 18c et 19c.

Oracle constitue un \textbf{écosystème logiciel complet}, comprenant notamment :

\begin{itemize}
    \item des serveurs de données et utilitaires (SQL*Plus, import/export loader, \dots) ;
    \item des environnements de développement (web, Java, pré-compilateurs, \dots) ;
    \item des outils d’aide à la décision (data warehouse, Oracle Discoverer, \dots) ;
    \item des progiciels client-serveur et Internet (Oracle Applications, Oracle e-Business Suite).
    \item $\cdots$
\end{itemize}

\section{Structure générale d’une requête}

Pour illustrer les requêtes SQL, on utilise une base de données exemple composée des tables suivantes :

\begin{itemize}
    \item \textbf{Client} (\underline{noClient}, nom, prénom, ddn, rue, CP, ville)
    \item \textbf{Produit} (\underline{noProduit}, libellé, prixUnitaire, \textit{noFournisseur})
    \item \textbf{Fournisseur} (\underline{noFournisseur}, raisonSociale)
    \item \textbf{Commande} (\underline{\textit{noClient}}, \underline{\textit{noProduit}}, \underline{dateCommande} , quantité)
\end{itemize}

Les attributs \underline{soulignés} désignent les \textbf{clés primaires} ;  
les attributs en \textit{italique} correspondent aux \textbf{clés étrangères}.

\subsection*{Conventions de notation}

\begin{itemize}
    \item \texttt{[élément]} : élément optionnel.
    \item \texttt{élément1 | élément2} : choisir l’un des deux éléments.
    \item \texttt{\{ élémentx | élémenty \}} : liste de choix possibles.
    \item \texttt{\{élémentx\} | \{élémenty\}} : choix entre deux ensembles.
\end{itemize}

\subsection{Syntaxe générale de la commande \texttt{SELECT}}

La commande \texttt{SELECT} permet d’interroger les données d’une ou plusieurs tables. Sa syntaxe générale est :

\begin{sqlcodebox}
\begin{lstlisting}[language=SQL]
SELECT * | {[ALL | DISTINCT]
        expression [[AS] nomColonne]
        [, expression [[AS] nomColonne]]... } 
FROM relation [alias] [, relation [alias] ...]
[WHERE condition]
[GROUP BY nomColonne [, nomColonne]...]
[HAVING condition]
[ORDER BY nomColonne [ASC | DESC]
         [, nomColonne [ASC | DESC] ...]];
\end{lstlisting}
\end{sqlcodebox}

Deux commandes \texttt{SELECT} peuvent être combinées :

\begin{sqlcodebox}
\begin{lstlisting}[language=SQL]
CdeSelect {UNION | INTERSECT | EXCEPT} CdeSelect
\end{lstlisting}
\end{sqlcodebox}

\subsection{Oracle-SQL et la casse des caractères}

Dans Oracle, il n’existe \textbf{aucune différence} entre identifiants écrits en majuscules ou en minuscules : les noms de tables, attributs ou contraintes sont insensibles à la casse.

La seule exception concerne la \textbf{comparaison de chaînes de caractères}, où \texttt{'Dupont'} et \texttt{'dupont'} sont considérés comme différents.

Exemple :

\begin{sqlcodebox}
\begin{lstlisting}[language=SQL]
SELECT *
FROM Client
WHERE nom = 'Dupont';
\end{lstlisting}
\end{sqlcodebox}

Les instructions suivantes sont équivalentes :

\begin{sqlcodebox}
\begin{lstlisting}[language=SQL]
SELECT *
FROM CLIENT
WHERE NOM = 'Dupont';
\end{lstlisting}
\end{sqlcodebox}







\section{Projection d'une relation}

La \textbf{projection} consiste à extraire seulement certains attributs d’une relation.  
Par exemple, pour récupérer les \textbf{numéros de clients} et les \textbf{dates de commande} de toutes les commandes, on utilise :

\begin{sqlcodebox}
\begin{lstlisting}[language=SQL]
SELECT noClient, dateCommande
FROM Commande;
\end{lstlisting}
\end{sqlcodebox}

Résultat :

\begin{tabular}{|c|c|}
    \hline
    \textbf{noClient} & \textbf{dateCommande}  \\
    \hline
    100 & 04/05/2003 \\
200 & 12/04/2003 \\
100 & 05/01/2004 \\
300 & 25/02/2004 \\
400 & 30/01/2004 \\
400 & 30/01/2004 \\
    \hline
\end{tabular}

Cette requête est équivalente à :

\begin{sqlcodebox}
\begin{lstlisting}[language=SQL]
SELECT ALL noClient, dateCommande
FROM Commande;
\end{lstlisting}
\end{sqlcodebox}

\subsection{La clause DISTINCT}

Pour éliminer les doublons, on utilise la clause \texttt{DISTINCT}.  
La requête suivante retourne chaque couple \texttt{(noClient, dateCommande)} une seule fois :

\begin{sqlcodebox}
\begin{lstlisting}[language=SQL]
SELECT DISTINCT noClient, dateCommande
FROM Commande;
\end{lstlisting}
\end{sqlcodebox}

Résultat :

\begin{tabular}{|c|c|}
    \hline
    \textbf{noClient} & \textbf{dateCommande}  \\
    \hline
    100 & 04/05/2003 \\
200 & 12/04/2003 \\
100 & 05/01/2004 \\
300 & 25/02/2004 \\
400 & 30/01/2004 \\
    \hline
\end{tabular}

Expression algébrique équivalente :

\[
\pi_{\text{noClient},\ \text{dateCommande}}(\text{Commande})
\]

%----------------------------------------------------

\section{Restriction d’une relation}

La \textbf{restriction} permet de sélectionner les tuples satisfaisant une condition logique.  
Exemple : on veut les articles dont le prix est inférieur à 20\,€ et dont le numéro est supérieur à 30 :

\begin{sqlcodebox}
\begin{lstlisting}[language=SQL]
SELECT *
FROM Article
WHERE prixUnitaire < 20 AND noArticle > 30;
\end{lstlisting}
\end{sqlcodebox}

Résultat :


\begin{tabular}{|c|c|c|}
    \hline
    \textbf{noArticle} & \textbf{libellé} & \textbf{prixUnitaire} \\
    \hline
31 & Brosse & 12 \\
40 & Tableau & 9.99 \\
50 & Visseuse & 19.99 \\
    \hline
\end{tabular}

Expression algébrique équivalente :

\[
\sigma_{\text{prixUnitaire} < 20\ \wedge\ \text{noArticle} > 30}(\text{Article})
\]

\subsection{Restriction : syntaxe d'une condition}

La syntaxe générale est :

\begin{sqlcodebox}
\begin{lstlisting}[language=SQL]
SELECT *
FROM nomRelation
WHERE condition;
\end{lstlisting}
\end{sqlcodebox}

Les formes possibles d’une condition :

\begin{itemize}
    \item \texttt{conditionSimple}
    \item \texttt{(condition)}
    \item \texttt{NOT (condition)}
    \item \texttt{condition AND condition}
    \item \texttt{condition OR condition}
\end{itemize}

Une \texttt{conditionSimple} peut être de la forme :

\begin{itemize}
    \item \texttt{expression =, <, >, <=, >=, <> expression}
    \item \texttt{expression [NOT] BETWEEN valeur1 AND valeur2}
    \item \texttt{expression IS [NOT] NULL}
    \item \texttt{expression [NOT] IN (liste)}
    \item \texttt{expression [NOT] LIKE patron}
\end{itemize}

\subsection{Restriction : exemples}

\textbf{Produits dont le prix est compris entre 50 et 100\,€}

\begin{sqlcodebox}
\begin{lstlisting}[language=SQL]
SELECT *
FROM Produit
WHERE prixUnitaire >= 50
  AND prixUnitaire <= 100;
\end{lstlisting}
\end{sqlcodebox}

\textbf{Produits dont le prix est inférieur à 50\,€ ou supérieur à 100\,€}

\begin{sqlcodebox}
\begin{lstlisting}[language=SQL]
SELECT *
FROM Produit
WHERE prixUnitaire < 50
   OR prixUnitaire > 100;
\end{lstlisting}
\end{sqlcodebox}

Écriture équivalente avec \texttt{BETWEEN} :

\begin{sqlcodebox}
\begin{lstlisting}[language=SQL]
SELECT *
FROM Produit
WHERE prixUnitaire BETWEEN 50 AND 100;
\end{lstlisting}
\end{sqlcodebox}

\subsection{Notion de valeur indéfinie (NULL)}

Une valeur \texttt{NULL} indique qu’un attribut est \textbf{inconnu} ou \textbf{indéfini}.  
Règles importantes :

\begin{itemize}
    \item On ne peut tester \texttt{NULL} qu’avec : \texttt{IS NULL} ou \texttt{IS NOT NULL}.
    \item Toute expression arithmétique impliquant \texttt{NULL} retourne \texttt{NULL}.
    \item Toute comparaison impliquant \texttt{NULL} retourne \texttt{UNKNOWN}.
\end{itemize}

\subsection{Restriction : opérateurs IS NULL et LIKE}

\textbf{Commandes dont la quantité est indéterminée}

\begin{sqlcodebox}
\begin{lstlisting}[language=SQL]
SELECT *
FROM Commande
WHERE quantite IS NULL;
\end{lstlisting}
\end{sqlcodebox}

\textbf{Clients dont le nom commence par B, finit par B et contient au moins 3 caractères}

\begin{sqlcodebox}
\begin{lstlisting}[language=SQL]
SELECT *
FROM Client
WHERE nom LIKE 'B_\%B';
\end{lstlisting}
\end{sqlcodebox}

Rappels sur les patrons \texttt{LIKE} :

\begin{itemize}
    \item \texttt{\%} représente une suite de 0 à $n$ caractères.
    \item \texttt{\_} représente exactement un caractère.
\end{itemize}

\subsection{Restriction : opérateur IN}

 \textbf{Clients dont le nom appartient à un ensemble donné}

\begin{sqlcodebox}
\begin{lstlisting}[language=SQL]
SELECT prenom
FROM Client
WHERE nom IN ('Dupond', 'Durant', 'Martin');
\end{lstlisting}
\end{sqlcodebox}

\textbf{Clients dont le nom n'appartient pas à cet ensemble}

\begin{sqlcodebox}
\begin{lstlisting}[language=SQL]
SELECT prenom
FROM Client
WHERE nom NOT IN ('Dupond', 'Durant', 'Martin');
\end{lstlisting}
\end{sqlcodebox}

%----------------------------------------------------

\section{Restriction et projection}

On peut combiner restriction et projection.  
Exemple : obtenir les numéros de clients et dates de commande pour les commandes postérieures au 01/01/2004 :

\begin{sqlcodebox}
\begin{lstlisting}[language=SQL]
SELECT noClient, dateCommande
FROM Commande
WHERE dateCommande > '01/01/2004';
\end{lstlisting}
\end{sqlcodebox}

Résultat :

\begin{tabular}{|c|c|}
    \hline
    \textbf{noClient} & \textbf{dateCommande}  \\
    \hline
    100 & 05/01/2004  \\
    300 & 25/02/2004  \\
    400 & 30/01/2004 \\
    \hline
\end{tabular}

Expression algébrique équivalente :

\[
\pi_{\text{noClient},\ \text{dateCommande}}
\bigl(
    \sigma_{\text{dateCommande} > '01/01/2004'}(\text{Commande})
\bigr)
\]



































\section{Jointure — Produit cartésien}

La jointure débute par l’opération la plus simple : le \textbf{produit cartésien}.  
Lorsqu’on liste plusieurs relations dans la clause \texttt{FROM}, sans condition supplémentaire, SQL génère toutes les combinaisons possibles de leurs tuples.

\begin{itemize}
    \item Exemple : produire toutes les paires possibles Client × Commande.
\end{itemize}

\begin{sqlcodebox}
\begin{lstlisting}[language=SQL]
SELECT *
FROM Client, Commande;
\end{lstlisting}
\end{sqlcodebox}

Expression algébrique équivalente :
\[
\text{Client} \times \text{Commande}
\]


\subsection{Jointure}

Une jointure SQL correspond en fait à la combinaison :
\[
\text{Produit cartésien} + \text{Restriction} + \text{Projection}
\]

Forme générale :

\begin{sqlcodebox}
\begin{lstlisting}[language=SQL]
SELECT attribut1 [, attribut2, ... ]
FROM relation1, relation2 [, relation3, ...]
WHERE condition;
\end{lstlisting}
\end{sqlcodebox}

Expression algébrique équivalente :
\[
\pi_{\text{attributs}} \big( \sigma_{\text{condition}} (\text{relation1} 
\times \text{relation2} \times \cdots ) \big)
\]


\subsection{Jointure : exemple de requête}

Afficher la liste des commandes accompagnées du nom du client :

\begin{sqlcodebox}
\begin{lstlisting}[language=SQL]
SELECT nom, Client.noClient,
       noProduit, dateCommande, quantite
FROM Commande, Client
WHERE Client.noClient = Commande.noClient;
\end{lstlisting}
\end{sqlcodebox}

\textbf{Remarque :} lorsqu’un attribut existe dans plusieurs relations, il faut le préfixer par le nom de la relation pour éviter toute ambiguïté.


\subsection{Jointure : autre exemple de requête}

Produire la liste des produits commandés en quantité supérieure à 100 et dont le prix dépasse 1000\,€.  
On affiche : numéro du produit, libellé, prix unitaire, date de commande.

\begin{sqlcodebox}
\begin{lstlisting}[language=SQL]
SELECT Produit.noProduit, libelle, prixUnitaire, date
FROM Produit, Commande
WHERE quantite > 100
  AND prixUnitaire > 1000
  AND Produit.noProduit = Commande.noProd;
\end{lstlisting}
\end{sqlcodebox}

\subsection{Jointure : utilisation d’alias}

Les alias permettent d’alléger l’écriture, surtout lorsqu’on joint plusieurs relations.

Exemple : afficher les commandes avec le nom et le numéro du client.

\begin{sqlcodebox}
\begin{lstlisting}[language=SQL]
SELECT C2.noClient, nom, date, quantite
FROM Commande C1, Client C2
WHERE C1.noClient = C2.noClient;
\end{lstlisting}
\end{sqlcodebox}

\begin{itemize}
    \item \texttt{Commande} est renommée \texttt{C1}
    \item \texttt{Client} est renommée \texttt{C2}
\end{itemize}


\subsection{Quelques règles de nommage}

\textbf{Nom d’une colonne dans le résultat :}

\begin{itemize}
    \item Par défaut : nom de l’attribut correspondant.
\end{itemize}

\begin{sqlcodebox}
\begin{lstlisting}[language=SQL]
SELECT * FROM Produit;
\end{lstlisting}
\end{sqlcodebox}

Résultat : \texttt{noProduit}, \texttt{libellé}, \texttt{prixUnitaire}

\medskip

\textbf{Renommage possible :}

\begin{sqlcodebox}
\begin{lstlisting}[language=SQL]
SELECT noProduit AS "Numero produit"
FROM Produit;
\end{lstlisting}
\end{sqlcodebox}

Résultat : colonne \texttt{Numéro produit}.


\section{Opérations ensemblistes (UNION, INTERSECT, EXCEPT)}

On peut combiner deux commandes \texttt{SELECT} avec les opérateurs :
\begin{itemize}
    \item \texttt{UNION} : union ensembliste
    \item \texttt{INTERSECT} : intersection
    \item \texttt{EXCEPT} : différence
\end{itemize}

Exemple : trouver les employés qui sont aussi passagers.

\begin{center}
\begin{tabular}{c c c c}
\multicolumn{2}{c}{\textbf{Employé}} & 
\multicolumn{2}{c}{\textbf{Passager}} \\
\hline
10 & Henry John & 4 & Harry Peter \\
15 & Conrad James & 78 & Conrad James \\
35 & Jenqua Jessica & 9 & Land Robert \\
46 & Leconte Jean & 466 & Leconte Jean 
\end{tabular}
\end{center}

Requête :

\begin{sqlcodebox}
\begin{lstlisting}[language=SQL]
(SELECT nomEmp AS nom, prenomEmp AS prenom
 FROM Employe)
INTERSECT
(SELECT nomPass AS nom, prenomPass AS prenom
 FROM Passager);
\end{lstlisting}
\end{sqlcodebox}


\section{Expression de calcul dans les colonnes projetées}

On peut intégrer des expressions arithmétiques directement dans la liste des résultats.

Exemple : afficher le prix TTC (taxe 15\%) :

\begin{sqlcodebox}
\begin{lstlisting}[language=SQL]
SELECT noArticle, prixUnitaire,
       prixUnitaire * 1.15 AS prixTTC
FROM Article;
\end{lstlisting}
\end{sqlcodebox}


\subsection{Expression de calcul dans la condition (WHERE)}

Les expressions peuvent aussi apparaître dans les conditions.

Produits dont le prix TTC dépasse 20\,€ :

\begin{sqlcodebox}
\begin{lstlisting}[language=SQL]
SELECT noArticle
FROM Article
WHERE prixUnitaire * 1.15 > 20;
\end{lstlisting}
\end{sqlcodebox}

Elles peuvent également utiliser des fonctions prédéfinies.

Exemple : commandes du jour :

\begin{sqlcodebox}
\begin{lstlisting}[language=SQL]
SELECT *
FROM Commande
WHERE dateCommande = CURRENT_DATE;
\end{lstlisting}
\end{sqlcodebox}


\section{Requêtes imbriquées}

Une requête peut apparaître à l’intérieur d’une autre, notamment dans la clause \texttt{WHERE}.  
On parle alors de \textbf{sous-requête} ou \textbf{requête imbriquée}.

Forme générale :

\begin{sqlcodebox}
\begin{lstlisting}[language=SQL]
SELECT colonne(s)
FROM relation(s)
WHERE expression [NOT] IN (sous-requete)
   OR EXISTS (sous-requete)
   OR NOT EXISTS (sous-requete);
\end{lstlisting}
\end{sqlcodebox}


\subsection{Requêtes imbriquées : IN / NOT IN}

Exemple : noms des clients ayant passé une commande le 24/10/2000.

\begin{sqlcodebox}
\begin{lstlisting}[language=SQL]
SELECT nom
FROM Client
WHERE noClient IN (SELECT noClient
                   FROM Commande
                   WHERE dateCommande = '24/10/2000');
\end{lstlisting}
\end{sqlcodebox}


\subsection{Requêtes imbriquées : EXISTS / NOT EXISTS}

Clients ayant passé au moins une commande :

\begin{sqlcodebox}
\begin{lstlisting}[language=SQL]
SELECT *
FROM Client C1
WHERE EXISTS (
  SELECT *
  FROM Commande C2
  WHERE C1.noClient = C2.noClient
);
\end{lstlisting}
\end{sqlcodebox}

Clients n’ayant passé aucune commande :

\begin{sqlcodebox}
\begin{lstlisting}[language=SQL]
SELECT *
FROM Client C1
WHERE NOT EXISTS (
  SELECT *
  FROM Commande C2
  WHERE C1.noClient = C2.noClient
);
\end{lstlisting}
\end{sqlcodebox}

\section{Fonctions d’agrégation}

Les fonctions d’agrégation permettent d’appliquer un calcul sur un ensemble de valeurs (souvent un groupe de lignes) pour produire une valeur unique.

La forme générale est :

\begin{sqlcodebox}
\begin{lstlisting}[language=SQL]
SELECT fctAgregation
FROM relation(s)
[WHERE condition];
\end{lstlisting}
\end{sqlcodebox}

Les principales fonctions sont :

\begin{itemize}
    \item \textbf{COUNT(*)} : nombre total de lignes du résultat.
    \item \textbf{COUNT([DISTINCT] expr)} : nombre de valeurs non NULL (distinctes éventuellement).
    \item \textbf{SUM(n)} : somme des valeurs de n (ignore les NULL).
    \item \textbf{MAX(n)} : valeur maximale.
    \item \textbf{MIN(n)} : valeur minimale.
    \item \textbf{AVG(n)} : moyenne des valeurs (ignore les NULL).
\end{itemize}

Ici, $n$ désigne une expression numérique et $expr$ toute expression valide.

\subsection*{Exemples}

\textbf{Nombre total d’articles et prix moyen :}

\begin{sqlcodebox}
\begin{lstlisting}[language=SQL]
SELECT COUNT(*) AS nbArticles,
       AVG(prixUnitaire) AS prixMoyen
FROM Article;
\end{lstlisting}
\end{sqlcodebox}

\textbf{Nombre de prix non NULL :}

\begin{sqlcodebox}
\begin{lstlisting}[language=SQL]
SELECT COUNT(prixUnitaire) AS nbPrixNonNull
FROM Article;
\end{lstlisting}
\end{sqlcodebox}

\textbf{Nombre de prix distincts :}

\begin{sqlcodebox}
\begin{lstlisting}[language=SQL]
SELECT COUNT(DISTINCT prixUnitaire) AS nbPrix
FROM Article;
\end{lstlisting}
\end{sqlcodebox}

% ---------------------------------------------------------

\section{Partition de relations (GROUP BY)}

La clause \texttt{GROUP BY} permet de regrouper les lignes selon une ou plusieurs colonnes, puis d’appliquer une fonction d’agrégation sur chaque groupe.

\subsection*{Exemple}

Nombre de produits commandés par client :

\begin{sqlcodebox}
\begin{lstlisting}[language=SQL]
SELECT noClient, COUNT(*) AS totalProduits
FROM Commande
GROUP BY noClient;
\end{lstlisting}
\end{sqlcodebox}

\textbf{Principe :}
\begin{enumerate}
    \item Les commandes sont regroupées selon \texttt{noClient}.
    \item Pour chaque groupe, on calcule le nombre d’éléments.
\end{enumerate}

Important : \textit{chaque expression présente dans le SELECT doit avoir une valeur unique par groupe.}

\subsection*{Exemple avec condition préalable}

Quantité totale commandée par client, hors produit \texttt{F565} :

\begin{sqlcodebox}
\begin{lstlisting}[language=SQL]
SELECT noClient, SUM(quantite)
FROM Commande
WHERE noProduit <> 'F565'
GROUP BY noClient;
\end{lstlisting}
\end{sqlcodebox}

Étapes :
\begin{enumerate}
    \item La clause WHERE élimine les commandes du produit F565.
    \item Les lignes restantes sont regroupées par client.
    \item Pour chaque groupe, la somme des quantités est calculée.
\end{enumerate}

% ---------------------------------------------------------

\subsection{Restriction des groupes : HAVING}

La clause \texttt{HAVING} agit comme un filtre, mais \textbf{sur les groupes}.  
Elle ne s’utilise qu’en présence d’un \texttt{GROUP BY}.

\subsection*{Exemple}

Quantité moyenne commandée par produit pour les produits ayant reçu plus de 3 commandes (en ignorant le client C47) :

\begin{sqlcodebox}
\begin{lstlisting}[language=SQL]
SELECT noProduit, AVG(quantite)
FROM Commande
WHERE noClient <> 'C47'
GROUP BY noProduit
HAVING COUNT(*) > 3;
\end{lstlisting}
\end{sqlcodebox}

Étapes :
\begin{enumerate}
    \item Élimination des commandes du client C47 (WHERE).
    \item Regroupement par produit.
    \item Élimination des groupes comptant moins de 3 commandes (HAVING).
    \item Calcul de la moyenne pour les groupes restants.
\end{enumerate}

% ---------------------------------------------------------

\subsection{Partition sur plusieurs colonnes}

Il est possible de grouper sur plusieurs attributs.

\subsection*{Exemple}

Nombre de produits commandés par client \textit{et par date} :

\begin{sqlcodebox}
\begin{lstlisting}[language=SQL]
SELECT noClient, dateCommande, COUNT(noProduit) AS nbProduits
FROM Commande
GROUP BY noClient, dateCommande;
\end{lstlisting}
\end{sqlcodebox}

Chaque combinaison \texttt{(noClient, dateCommande)} forme un groupe distinct.

% ---------------------------------------------------------

\subsection{Sous-requête et fonctions d’agrégation}

Une fonction d’agrégation ne peut pas être utilisée directement dans un SELECT si elle produit un résultat incompatible avec les autres colonnes.

\textbf{Exemple incorrect :}

\begin{sqlcodebox}
\begin{lstlisting}[language=SQL]
SELECT noProduit, MAX(prixUnitaire)
FROM Produit;
\end{lstlisting}
\end{sqlcodebox}

Cette requête est invalide car \texttt{noProduit} n’est pas unique.

\medskip

Pour contourner, on utilise une sous-requête :

\begin{sqlcodebox}
\begin{lstlisting}[language=SQL]
SELECT noProduit, libelle
FROM Produit
WHERE prixUnitaire > (
    SELECT AVG(prixUnitaire)
    FROM Produit
);
\end{lstlisting}
\end{sqlcodebox}

% ---------------------------------------------------------

\section{Tri du résultat d’une requête (ORDER BY)}

La clause \texttt{ORDER BY} permet de trier le résultat d'une requête selon une ou plusieurs colonnes.

\begin{itemize}
    \item \textbf{ASC} : ordre croissant (défaut)
    \item \textbf{DESC} : ordre décroissant
\end{itemize}

\subsection*{Exemple}

Liste des commandes triées par numéro de client (croissant), puis par date (décroissant) :

\begin{sqlcodebox}
\begin{lstlisting}[language=SQL]
SELECT *
FROM Commande
ORDER BY noClient, dateCommande DESC;
\end{lstlisting}
\end{sqlcodebox}

\section{Mise à jour des données}

Les opérations de mise à jour permettent d’ajouter, de modifier ou de supprimer
des tuples dans une relation. Elles s'effectuent à l’aide des instructions
\texttt{INSERT}, \texttt{UPDATE} et \texttt{DELETE}.

\subsection{Ajout de tuples (INSERT)}

L’instruction \texttt{INSERT} permet d’insérer un ou plusieurs tuples dans une relation.
Sa forme générale est la suivante :

\begin{sqlcodebox}
\begin{lstlisting}[language=SQL]
INSERT INTO nom_relation [(attribut, attribut, ...)]
VALUES (valeur, valeur, ...)
\end{lstlisting}
\end{sqlcodebox}

ou bien :

\begin{sqlcodebox}
\begin{lstlisting}[language=SQL]
INSERT INTO nom_relation [(attribut, attribut, ...)]
commande-SELECT
\end{lstlisting}
\end{sqlcodebox}

\textbf{Exemple : ajout d'un produit}

\begin{sqlcodebox}
\begin{lstlisting}[language=SQL]
INSERT INTO Produit VALUES (430, 'lecteur DVD', 9.99);
\end{lstlisting}
\end{sqlcodebox}

Tous les attributs doivent être fournis dans le même ordre que dans la définition
de la table.

\textbf{Exemple : insertion provenant d’un SELECT}

\begin{sqlcodebox}
\begin{lstlisting}[language=SQL]
INSERT INTO Commande (noClient, noProduit)
SELECT noClient, noProduit FROM Produit, Client;
\end{lstlisting}
\end{sqlcodebox}

Les attributs non mentionnés dans la liste sont automatiquement positionnés à
\texttt{NULL}.

\subsection{Modification de tuples (UPDATE)}

L’instruction \texttt{UPDATE} permet de modifier la valeur d’un ou plusieurs attributs
pour les tuples vérifiant une condition.

\begin{sqlcodebox}
\begin{lstlisting}[language=SQL]
UPDATE nom_relation
SET attribut = expression
    [, attribut = expression ...]
[WHERE condition];
\end{lstlisting}
\end{sqlcodebox}

\textbf{Exemple : modifier le nom du client n°3}

\begin{sqlcodebox}
\begin{lstlisting}[language=SQL]
UPDATE Client
SET nom = 'Durand'
WHERE noClient = 3;
\end{lstlisting}
\end{sqlcodebox}

\textbf{Exemple : augmenter de 5\% le prix de certains produits}

\begin{sqlcodebox}
\begin{lstlisting}[language=SQL]
UPDATE Produit
SET prixUnitaire = prixUnitaire * 1.05
WHERE libelle IN ('CD-ROM', 'DVD', 'ZIP');
\end{lstlisting}
\end{sqlcodebox}

\paragraph{Exemple : mise à jour basée sur une sous-requête}

Augmenter de 10\% les prix des produits fournis par le fournisseur « Titu » :

\begin{sqlcodebox}
\begin{lstlisting}[language=SQL]
UPDATE Produit p
SET prixUnitaire = prixUnitaire * 1.1
WHERE p.noFournisseur IN (
    SELECT f.noFournisseur
    FROM Fournisseur f
    WHERE f.raisonSociale = 'Titu'
);
\end{lstlisting}
\end{sqlcodebox}

\paragraph{Autre exemple}

Positionner le libellé du produit n°99 et lui attribuer le prix moyen :

\begin{sqlcodebox}
\begin{lstlisting}[language=SQL]
UPDATE Produit
SET libell" = "produiTest",
    prixUnitaire = (SELECT AVG(prixUnitaire) FROM Produit)
WHERE noProduit = 99;
\end{lstlisting}
\end{sqlcodebox}

\subsection{Suppression de tuples (DELETE)}

L’instruction \texttt{DELETE} supprime les tuples correspondant à une condition.

\begin{sqlcodebox}
\begin{lstlisting}[language=SQL]
DELETE FROM relation
[WHERE condition];
\end{lstlisting}
\end{sqlcodebox}

\textbf{Exemple : supprimer les clients habitant Metz}

\begin{sqlcodebox}
\begin{lstlisting}[language=SQL]
DELETE FROM Client
WHERE ville = 'Metz';
\end{lstlisting}
\end{sqlcodebox}

\textbf{Exemple : suppression totale}

\begin{sqlcodebox}
\begin{lstlisting}[language=SQL]
DELETE FROM Client;
\end{lstlisting}
\end{sqlcodebox}

Toutes les lignes sont supprimées, mais le schéma de la table demeure.

% ---------------------------------------------------------------------------

\section{CREATE TABLE}

L’instruction \texttt{CREATE TABLE} permet de créer une nouvelle table en précisant
ses attributs et ses contraintes.

\begin{sqlcodebox}
\begin{lstlisting}[language=SQL]
CREATE TABLE nomTable (
    definitionAttribut,
    ...
    definitionContrainte
);
\end{lstlisting}
\end{sqlcodebox}

\subsection*{Syntaxe de définition d’un attribut}

\begin{sqlcodebox}
\begin{lstlisting}[language=SQL]
nomAttribut type
    [DEFAULT valeur]
    [NULL | NOT NULL]
    [UNIQUE | PRIMARY KEY]
    [REFERENCES autreTable(attribut)]
    [CONSTRAINT nom CHECK (condition)]
\end{lstlisting}
\end{sqlcodebox}

Il n’y a pas de distinction entre minuscules et majuscules dans les noms.

\subsection*{Syntaxe de définition d’une contrainte}

\begin{sqlcodebox}
\begin{lstlisting}[language=SQL]
[CONSTRAINT nom]
PRIMARY KEY (liste_attributs)
UNIQUE (liste_attributs)
FOREIGN KEY (liste_attributs)
    REFERENCES tableBis (liste_attributs)
    [ON DELETE regle]
    [ON UPDATE regle]
CHECK (condition)
\end{lstlisting}
\end{sqlcodebox}

% ---------------------------------------------------------------------------

\subsection{Types de données dans ORACLE}

\paragraph{Types numériques}

Les types ANSI INT, SMALLINT, FLOAT, REAL… sont convertis en type ORACLE
\texttt{NUMBER}.

\begin{center}
\begin{tabular}{ll}
\textbf{Type SQL ANSI} & \textbf{Type ORACLE} \\
NUMERIC(p,c), DECIMAL(p,c) & NUMBER(p,c) \\
INTEGER, SMALLINT & NUMBER(38) \\
FLOAT(b), REAL & NUMBER \\
\end{tabular}
\end{center}

\paragraph{Chaînes de caractères}

\begin{itemize}
\item \texttt{CHAR(n)} : taille fixe ($n \leq 2000$)
\item \texttt{VARCHAR(n)} : taille variable ($n \leq 4000$)
\end{itemize}

% ---------------------------------------------------------------------------

\subsection{Contraintes sur le domaine d’un attribut}

\paragraph{Contrainte NOT NULL}

L’utilisateur doit fournir une valeur.

\begin{sqlcodebox}
\begin{lstlisting}[language=SQL]
CREATE TABLE Client (
    noClient INTEGER NOT NULL,
    nom VARCHAR(50) NOT NULL,
    ddn VARCHAR(15),
    telephone VARCHAR(15)
);
\end{lstlisting}
\end{sqlcodebox}

\textbf{Contrainte CHECK sur un attribut}

\begin{sqlcodebox}
\begin{lstlisting}[language=SQL]
CREATE TABLE Client (
    noClient INTEGER NOT NULL
        CHECK (noClient > 0 AND noClient < 1000),
    nom VARCHAR(50) NOT NULL,
    ddn VARCHAR(15)
);
\end{lstlisting}
\end{sqlcodebox}

\textbf{Contrainte CHECK sur plusieurs attributs}

\begin{sqlcodebox}
\begin{lstlisting}[language=SQL]
CREATE TABLE Produit (
    noProduit INTEGER NOT NULL,
    libelle VARCHAR(50),
    prixUnitaire NUMBER(10,2) NOT NULL,
    CHECK (noProduit <= 100 OR prixUnitaire > 15)
);
\end{lstlisting}
\end{sqlcodebox}

\textbf{Valeur par défaut}

\begin{sqlcodebox}
\begin{lstlisting}[language=SQL]
CREATE TABLE employe (
    numEmp INTEGER,
    numTel VARCHAR(15) DEFAULT 'confidentiel' NOT NULL
);
\end{lstlisting}
\end{sqlcodebox}

% ---------------------------------------------------------------------------

\subsection{Contrainte de clé primaire (PRIMARY KEY)}

La clé primaire identifie de manière unique chaque tuple et ne peut contenir
de valeur \texttt{NULL}.

\begin{sqlcodebox}
\begin{lstlisting}[language=SQL]
CREATE TABLE Client (
    noClient INTEGER PRIMARY KEY,
    nom VARCHAR(50) NOT NULL,
    ddn VARCHAR(15)
);
\end{lstlisting}
\end{sqlcodebox}

\textbf{Clé primaire composée}

\begin{sqlcodebox}
\begin{lstlisting}[language=SQL]
CREATE TABLE Commande (
    noClient INTEGER,
    noProduit INTEGER,
    dateCommande DATE,
    quantite INTEGER,
    PRIMARY KEY (noClient, noProduit, dateCommande)
);
\end{lstlisting}
\end{sqlcodebox}

\textbf{Contrainte d’unicité (UNIQUE)}

\begin{sqlcodebox}
\begin{lstlisting}[language=SQL]
CREATE TABLE Citoyen (
    numSecu INTEGER PRIMARY KEY,
    nom VARCHAR(50),
    prenom VARCHAR(50),
    ddn DATE,
    noPassport INTEGER UNIQUE
);
\end{lstlisting}
\end{sqlcodebox}

\textbf{Contrainte d’intégrité référentielle (FOREIGN KEY)}

\begin{sqlcodebox}
\begin{lstlisting}[language=SQL]
CREATE TABLE Commande (
    noCommande INTEGER PRIMARY KEY,
    noClient INTEGER REFERENCES Client(noClient),
    noProduit INTEGER,
    dateCommande DATE,
    quantite INTEGER,
    FOREIGN KEY (noProduit)
        REFERENCES Produit(noProduit)
);
\end{lstlisting}
\end{sqlcodebox}

Les tables référencées doivent exister et posséder une clé primaire ou une clé
unique.

% ---------------------------------------------------------------------------

\subsection{Création de schéma + instanciation}

\begin{sqlcodebox}
\begin{lstlisting}[language=SQL]
CREATE TABLE nomTable [(attributs)]
AS SELECT ...
\end{lstlisting}
\end{sqlcodebox}

Cette syntaxe :

\begin{enumerate}
\item crée la structure de la table,
\item insère les résultats d’un \texttt{SELECT}.
\end{enumerate}

\textbf{Exemple :}

\begin{sqlcodebox}
\begin{lstlisting}[language=SQL]
CREATE TABLE Client54 AS
SELECT noClient, nom
FROM Client
WHERE CP BETWEEN 54000 AND 54999;
\end{lstlisting}
\end{sqlcodebox}
